\scanpage{85}
\mspage{538}

\begin{center}
TABLE DES MATIERES
\end{center}

\noindent AVANT-PROPOS \dotfill

\noindent INTRODUCTION\dotfill

\noindent CHAPITRE \unsure{O|0} .- Préliminaires\dotfill

\noindent\quad \S\ 1. Foncteurs représentables\dotfill

\noindent\qquad 1.1. Foncteurs représentables\dotfill

\noindent\qquad 1.2. Produits fibrés dans une catégorie\dotfill

\noindent\qquad 1.3. Changement de base\dotfill

\noindent\qquad 1.4. Noyaux ; morphisme diagonal\dotfill

\noindent\qquad 1.5. Foncteurs adjoints\dotfill

\noindent\qquad 1.6. Structures algébriques dans les catégories\dotfill

\noindent\qquad 1.7. Morphismes fonctoriels représentables\dotfill

\noindent\quad \S\ 2. Compléments de topologie\dotfill

\noindent\qquad 2.1. Espaces irréductibles\dotfill

\noindent\qquad 2.2. Espaces localement noethériens\dotfill

\noindent\qquad 2.3. Ensembles quasi-constructibles et ensembles constructibles\dotfill

\margin{\handwritten{?}}
\noindent\qquad \deleted{2.4. Fonctions constructibles}\deleted{\mbox{............................}}

\noindent\qquad \deleted{2.5.}

\noindent\qquad 2.4. Ensembles constructibles dans les espaces noethériens

\noindent\qquad 2.5? Fonctions constructibles\dotfill

\noindent\qquad 2.6. Parties très denses d'un espace topologique\dotfill

\noindent\qquad 2.7. Quasi-homéomorphismes\dotfill

\noindent\qquad 2.8. Espaces de Jacobson\dotfill

\noindent\qquad 2.9. Espace sobre associé à un espace topologique\dotfill

\noindent\qquad 2.10.Applications ouvertes en un point \dotfill

\noindent\quad \S\ 3. Compléments sur les faisceaux\dotfill

\noindent\qquad 3.1. Faisceaux à valeurs dans une catégorie\dotfill

\noindent\qquad 3.2. Préfaisceaux sur une base d'ouverts\dotfill

\noindent\qquad 3.3. Recollement de faisceaux\dotfill

\noindent\qquad 3.4. Images directes de préfaisceaux\dotfill

\noindent\qquad 3.5. Images réci\deleted{\unsure{r}}p\deleted{p}roques de préfaisceaux\dotfill

\noindent\qquad 3.6. Faisceaux simples et faisceaux localement simples\dotfill

\noindent\qquad 3.7. Images réciproques de faisceaux de groupes ou d'anneaux\dotfill

% archive sheet number: 84
